\section{Self-Attention: Query, Key, Value e interacción directa}

La \term{self-attention} (autoatención) hace que la secuencia de entrada produzca al mismo tiempo query, key y value. Cada posición compara su query con todas las key y luego combina los value según los pesos obtenidos. Su diferencia con la encoder-decoder attention no está en que la fórmula sea completamente distinta, sino en que query, key y value provienen de la misma secuencia, por lo que los token pueden establecer directamente dependencias internas.

\subsection{Scaled dot-product attention}

Sea la matriz de entrada $X\in\mathbb{R}^{n\times d_{model}}$; mediante proyecciones lineales se obtiene
\[
Q=XW_Q,\qquad K=XW_K,\qquad V=XW_V.
\]
La scaled dot-product attention es
\[
\operatorname{Attention}(Q,K,V)
=\operatorname{softmax}\!\left(\frac{QK^\top}{\sqrt{d_k}}+M\right)V.
\]
$n$ es el número de token, $d_k$ es la dimensión de key/query, $QK^\top$ da todos los score query-key, $\sqrt{d_k}$ evita que el dot product crezca demasiado cuando aumenta la dimensión, y $M$ es una \term{mask} (máscara) opcional que suma $-\infty$ a las posiciones a las que no se permite acceder, de modo que, después del softmax, su peso es prácticamente cero.

\lecturefigure{slide-11-self-attention-formulation.jpg}{La self-attention produce query, key y value a partir de la misma secuencia.}{Video oficial, 08:00; Vaswani et al. 2017.}

\begin{knowledgebox}{Lectura de la figura: un token cumple tres papeles a la vez}
La query expresa «qué necesito ahora», la key expresa «con qué índice puedo ser encontrado» y el value expresa «qué contenido aporto si me eligen». Las tres provienen de proyecciones lineales distintas, de modo que una misma palabra puede usar un subespacio para determinar la relevancia y otro para transmitir información. La salida sigue siendo un conjunto de token representation, no un único vector en el que se comprime toda la secuencia.
\end{knowledgebox}

\begin{importantbox}{Intuición con un cálculo a mano de tres token}
Supongamos que los puntajes escalados entre una query y tres key son $[2,1,0]$; después del softmax quedan aproximadamente en $[0.665,0.245,0.090]$. Si los values son $v_1,v_2,v_3$, la salida es $0.665v_1+0.245v_2+0.090v_3$. Esto no equivale a recuperar un registro de una base de datos, sino a una mezcla ponderada diferenciable; el entrenamiento aprende simultáneamente el score space y el value content.
\end{importantbox}

\begin{lstlisting}[caption={Pseudocódigo mínimo de la scaled dot-product attention}]
scores = (query @ key.transpose(-2, -1)) / sqrt(key_dim)
scores = scores + mask
weights = softmax(scores, dim=-1)
output = weights @ value
\end{lstlisting}

\subsection{Rutas de información y paralelismo}

En una RNN, para que el primer token influya en el último normalmente hay que pasar por $O(n)$ pasos recurrentes; en una sola capa de full self-attention, dos posiciones cualesquiera interactúan directamente a través de una attention edge. Durante el entrenamiento, todas las query pueden calcularse en paralelo mediante multiplicación de matrices, lo cual se ajusta al alto rendimiento de los accelerator. El costo es que la score matrix explícita tiene $n\times n$ entradas.

\lecturefigure{slide-12-self-attention-diagram.jpg}{La self-attention despliega las relaciones token-to-token como conexiones directas que pueden calcularse en paralelo.}{Video oficial, 10:30; animación de clase, final state.}

\begin{knowledgebox}{Lectura de la figura: observe las conexiones, no solo los recuadros}
En la figura, los token de entrada pasan por proyecciones lineales que producen varios grupos de representaciones, y cada posición de salida agrega información proveniente de las demás posiciones. En comparación con una recurrent chain, la ruta de dependencia es más corta; en comparación con la convolution, las conexiones no son un kernel local fijo, sino que las determina dinámicamente el contenido. Lo que la figura no muestra es la división de trabajo entre attention head, la residual path ni el apilamiento de capas; por lo tanto, explica el routing, pero no representa un Transformer block completo.
\end{knowledgebox}

\begin{warningbox}{Una conexión directa no implica aprender automáticamente las relaciones correctas}
La self-attention permite que cualquier token interactúe con cualquier otro, pero el modelo sigue necesitando datos, objective, position, optimization y suficientes capas para aprender patrones útiles. Una capacidad alta también puede ajustarse a shortcut, spurious correlation o filtración de información; la accesibilidad es una condición para la capacidad, no una garantía de corrección.
\end{warningbox}

\subsection{Multi-head attention}

La \term{multi-head attention} (atención multicabeza) ejecuta en paralelo varias attention de menor dimensión:
\[
\operatorname{head}_r=\operatorname{Attention}(QW_r^Q,KW_r^K,VW_r^V),
\qquad
\operatorname{MHA}=\operatorname{Concat}(\operatorname{head}_1,\ldots,\operatorname{head}_H)W^O.
\]
Distintas head pueden aprender patrones locales, de largo alcance, sintácticos o posicionales, pero no se garantiza que cada head tenga un significado claro para una persona. La función principal de las múltiples cabezas es proporcionar varios subespacios de representación y varios routing channel en paralelo.

\begin{knowledgebox}{Términos clave: Q/K/V, head y mask}
\begin{tabularx}{\textwidth}{L{0.18\textwidth} X}
\toprule
Término & Problema que resuelve y relación con el curso \\
\midrule
Query & Expresa la necesidad de lectura de la posición actual; determina a qué key interroga. \\
Key & Proporciona el índice de coincidencia de las posiciones candidatas; junto con la query produce el score. \\
Value & Contenido que se mezcla según los pesos; la relevancia y el contenido transmitido pueden usar subespacios distintos. \\
Head & Proyección y attention channel independientes; amplían los tipos de relación y la capacidad de representación. \\
Mask & Impone límites a la información; la padding mask ignora las posiciones vacías y la causal mask impide ver el futuro. \\
\bottomrule
\end{tabularx}
\end{knowledgebox}

\subsection{Resumen de la sección}

El núcleo de la self-attention es la coincidencia aprendible entre query y key y la mezcla de value. Acorta las rutas de dependencia y aumenta el paralelismo del entrenamiento, pero genera una score matrix cuadrática y requiere position, mask, residual y normalization para convertirse en una arquitectura estable.
